\chapter{El \emph{planner}}
\label{cap:planner}

\begin{intuicion}
GraphHunter tiene cuatro operadores de matching ---DFS legacy con
prunes, LFTJ-K\(_4\), Yannakakis chain, Yannakakis tree-walk---.
El \emph{planner} es la pieza chiquita y fundamental que decide
cuál se usa para cada hunt. Inspecciona la \texttt{Hypothesis},
clasifica su shape y propiedades estructurales, y devuelve un
\texttt{DispatchPlan}. La política es deliberadamente simple: tres
reglas en orden, primer match gana, fallback al DFS. Este capítulo
documenta las tres reglas, justifica su orden, y explica cómo
agregar reglas nuevas sin romper backwards compat.
\end{intuicion}

\section{El \texttt{DispatchPlan}}

\begin{lstlisting}[language=Rust]
pub enum DispatchPlan {
    DfsMatch,    // default, generic; implementa cualquier hipótesis
    K4Lftj,      // M1/S3, sólo K_4 homogéneo
    Yannakakis,  // L1, $\alpha$-acíclica + flag GRAPHHUNTER_YANNAKAKIS
}
\end{lstlisting}

\noindent
Tres variantes, no más. \texttt{DfsMatch} es el fallback genérico.
\texttt{K4Lftj} y \texttt{Yannakakis} son los \emph{lifts}
especializados. La filosofía: \emph{conservadora}; sólo dispatchamos
a un especializado si \emph{sabemos} que la estructura lo permite.

\section{Las tres reglas}

\begin{lstlisting}[language=Rust]
pub fn plan(hypothesis: &Hypothesis) -> DispatchPlan {
    if name_marker_dispatch(hypothesis)   { return DispatchPlan::K4Lftj; }
    if structural_dispatch(hypothesis)    { return DispatchPlan::K4Lftj; }
    if yannakakis_dispatch(hypothesis)    { return DispatchPlan::Yannakakis; }
    DispatchPlan::DfsMatch
}
\end{lstlisting}

\noindent
Cada función es un test booleano sobre la \texttt{Hypothesis}. La
primera que retorna \texttt{true} gana.

\subsection{Regla 1: \texttt{name\_marker\_dispatch} (back-compat S3)}

\begin{lstlisting}[language=Rust]
fn name_marker_dispatch(hyp: &Hypothesis) -> bool {
    std::env::var("GRAPHHUNTER_K4_LFTJ").map_or(false, |v| v == "1")
        && hyp.name.to_lowercase().contains("k4-clique")
}
\end{lstlisting}

\noindent
Es la primera versión histórica del dispatch a LFTJ: el analista
nombra explícitamente su hipótesis con la string ``k4-clique'' y el
motor asume que es \(K_4\). Frágil, pero útil para debugging y
rollback. Hoy se mantiene para que tests viejos no se rompan; los
hunts nuevos no necesitan el marker.

\subsection{Regla 2: \texttt{structural\_dispatch} (M1)}

\begin{lstlisting}[language=Rust]
fn structural_dispatch(hyp: &Hypothesis) -> bool {
    let lftj_planner = std::env::var("GRAPHHUNTER_LFTJ_PLANNER")
        .map_or(true, |v| v == "1");  // default ON
    if !lftj_planner { return false; }

    // Detectar K_4 homogéneo:
    // - exactamente 6 pasos
    // - todos los pasos tienen el mismo EntityType en source y dest
    // - todos los pasos tienen el mismo RelationType
    // - sin predicados sobre metadata
    if hyp.steps.len() != 6 { return false; }
    let first = &hyp.steps[0];
    let same_entity_type = hyp.steps.iter().all(|s| {
        s.source_type == first.source_type && s.dest_type == first.dest_type
    });
    let same_rel_type = hyp.steps.iter().all(|s| s.rel_type == first.rel_type);
    let no_predicates = hyp.steps.iter().all(|s| s.predicates.is_empty());

    same_entity_type && same_rel_type && no_predicates
        && first.source_type == first.dest_type  // K_4 es homogéneo
}
\end{lstlisting}

\noindent
M1 (entregado abril 2026) destrabó el dispatch automático sin
requerir name marker. Inspecciona la estructura: 6 pasos
homogéneos, sin predicados, mismo tipo en todo $\Rightarrow$ es
\(K_4\) homogéneo $\Rightarrow$ LFTJ funciona.

\subsection{Regla 3: \texttt{yannakakis\_dispatch} (L1)}

\begin{lstlisting}[language=Rust]
fn yannakakis_dispatch(hyp: &Hypothesis) -> bool {
    let yannakakis_on = std::env::var("GRAPHHUNTER_YANNAKAKIS")
        .map_or(true, |v| v == "1");  // default ON
    if !yannakakis_on { return false; }

    // Cualquier hipótesis $\alpha$-acíclica aplica.
    // Tanto shape Chain como shape Graph dispatchan al matcher
    // de Yannakakis correspondiente.
    is_alpha_acyclic(hyp)
}
\end{lstlisting}

\noindent
La regla es la más general: \emph{cualquier hipótesis
\(\alpha\)-acíclica}. La función \texttt{is\_alpha\_acyclic} corre
GYO sobre el hipergrafo de la hipótesis (capítulo
\ref{cap:hipergrafos}) ---trivialmente rápido para hipótesis con
\(\leq 6\) pasos.

Notar que cubre las dos shapes:
\begin{itemize}
    \item \texttt{HypothesisShape::Chain}: dispatcha a Yannakakis
    chain (capítulo~\ref{cap:yannakakis-chain}).
    \item \texttt{HypothesisShape::Graph}: dispatcha a Yannakakis
    tree-walk (capítulo~\ref{cap:yannakakis-tree}).
\end{itemize}

\noindent
Adentro del matcher Yannakakis, otro \emph{branch} sobre
\texttt{shape} elige el sub-operador.

\section{¿Por qué este orden?}

Las tres reglas se chequean en orden. Razones:

\begin{enumerate}
    \item \textbf{name\_marker antes de structural}: el name marker
    permite force-override en tests y A/B comparisons. Si está
    seteado, gana sin importar la estructura.
    \item \textbf{structural antes de yannakakis}: \(K_4\) es
    cíclico, así que \texttt{yannakakis\_dispatch} lo rechazaría.
    Pero LFTJ-K\(_4\) le da mejor performance que el DFS para
    \(K_4\). Si chequeáramos Yannakakis primero, \(K_4\) caería al
    DFS innecesariamente.

    Además, la jerarquía estructural es estricta: \(K_4\)
    homogéneo (\texttt{structural}) es un subset de los hunts; los
    \(\alpha\)-acíclicos son otro subset; ambos son disjoint
    (\(K_4\) NO es \(\alpha\)-acíclico). Pero el orden importa
    sólo si una hipótesis cae en ambos ---imposible por la
    aciclicidad---. La razón real para el orden es de
    \emph{previsibilidad}: el lector que mira \texttt{plan()}
    sabe qué se intenta primero.
    \item \textbf{Fallback al DFS}: el DFS implementa cualquier
    hipótesis (incluso las que ningún operador especializado
    soporta). Es la red de seguridad.
\end{enumerate}

\section{Cuándo el planner ``se equivoca''}

El planner es \emph{estructural}: mira la shape, no el grafo. Es
posible que dispatchar a Yannakakis (\(O(|\E| + |Q_{\text{out}}|)\))
sea peor que DFS para una hipótesis donde NLF + k-hop ya filtraron
hasta dejar 10 starts válidos:

\begin{itemize}
    \item DFS desde 10 starts: $\sim$10~ms.
    \item Yannakakis full reducer sobre $|\E| = 10^7$: $\sim$300~ms.
\end{itemize}

\noindent
Hoy el planner no tiene heurísticos sobre el \emph{grafo}, sólo sobre
la \emph{hipótesis}. La razón es simplicidad: agregar costos
estimados requiere statistics actualizadas (cardinalidades,
selectividades) que el motor no mantiene. Para casos donde el
analista quiere forzar DFS, puede setear
\texttt{GRAPHHUNTER\_YANNAKAKIS=0}.

\section{Surface en UI: ``Perf: M+L1''}

El planner expone su elección al frontend vía \texttt{cmd\_get\_perf\_status}:

\begin{lstlisting}[language=Rust]
pub struct PerfStatus {
    pub lftj_planner: bool,        // M1
    pub k4_lftj_marker: bool,      // S3
    pub yannakakis: bool,           // L1
    pub plug_and_play_active: bool,
}
\end{lstlisting}

\noindent
El frontend renderiza una pildora ``Perf: M+L1'' (verde) cuando los
tres flags están on (default). Si el usuario forzó alguno a 0
(legacy path para debugging), la pildora se vuelve gris ``Perf:
legacy'' (capítulo~\ref{cap:plug-and-play}).

\section{Agregar una regla nueva: el procedimiento}

Cuando el equipo decide agregar un operador nuevo (digamos, LFTJ
para triángulo):

\begin{enumerate}
    \item Implementar el operador en
    \filepath{core/graph-engine/src/lftj\_triangle.rs}.
    \item Agregar variante a \texttt{DispatchPlan}:
    \texttt{TriangleLftj}.
    \item Escribir una función \texttt{triangle\_dispatch} con la
    detección estructural (3 pasos, ciclo, mismo tipo).
    \item Agregar a la cadena de \texttt{plan()}, antes de
    \texttt{yannakakis\_dispatch} (porque triángulo es cíclico,
    Yannakakis lo rechazaría):
    \begin{lstlisting}[language=Rust]
        if triangle_dispatch(hyp) { return DispatchPlan::TriangleLftj; }
    \end{lstlisting}
    \item Agregar el flag a \texttt{PerfStatus} y a la pildora.
    \item Default-on al startup (\filepath{apps/tauri/src-tauri/src/lib.rs})
    siguiendo la regla plug-and-play (capítulo~\ref{cap:plug-and-play}).
\end{enumerate}

\noindent
La política: cada nuevo operador es una rama nueva en
\texttt{plan()} con su detección estructural propia, default-on,
expuesto en UI.

\section{Resumen}

\begin{itemize}
    \item Planner: tres reglas en orden
    (\texttt{name\_marker}, \texttt{structural}, \texttt{yannakakis}),
    primero que matchea gana, fallback al DFS.
    \item \texttt{name\_marker}: legacy hardcoded
    para back-compat.
    \item \texttt{structural}: detecta \(K_4\) homogéneo via
    inspección de pasos.
    \item \texttt{yannakakis}: cualquier hipótesis
    \(\alpha\)-acíclica (chain o graph).
    \item DFS legacy es el fallback genérico; implementa cualquier
    hipótesis.
    \item El planner es estructural, no statistic-aware; cuando se
    equivoca, el usuario puede forzar via flags.
\end{itemize}

\begin{ejercicios}
\begin{enumerate}
    \item \starred{} Para cada una de estas hipótesis, identificá
    qué operador dispatcha el planner:
    \begin{itemize}
        \item \texttt{User -[Auth]-> Host -[Connect]-> IP}
        \item \texttt{c -[Knows]-> p1; c -[Knows]-> p2; p1 -[Knows]-> p2}
        \item Hipótesis con name=``k4-clique'' y 6 pasos heterogéneos
    \end{itemize}

    \item ¿Por qué \texttt{name\_marker\_dispatch} se chequea
    primero? Discutí: ¿qué se rompe si lo movemos al final?

    \item \starred{} Diseñá un detector estructural para la
    \emph{estrella de \(k\) brazos}. ¿Cómo se vería la función
    booleana?

    \item Si una hipótesis es \(\alpha\)-acíclica pero el usuario
    forzó \texttt{GRAPHHUNTER\_YANNAKAKIS=0}, ¿a qué operador cae?
    Justificá por qué la cadena de fallbacks es ``correcta''.

    \item Imaginá que el motor empieza a recolectar estadísticas
    de cardinalidades. ¿Cómo cambiarías \texttt{plan()} para
    aprovecharlas? Discutí pros y contras de un planner
    cost-based.
\end{enumerate}
\end{ejercicios}
